\section{Methodology: Rademacher-Bounded Polynomial Merging}
\label{sec:method}

We present \textbf{Rademacher-Bounded Polynomial Merging (RBPM)}, a mathematically rigorous framework designed to solve the overparameterization and transductive overfitting issues of post-hoc adaptive model merging. We treat the ensembling coefficients across network layers as a continuous global trajectory, deriving tight Rademacher complexity bounds and formulating a theoretically sound, regularized objective.

\subsection{Parameter Space Formulation}
Let $W_0^{(l)} \in \mathbb{R}^{D_l}$ represent the parameter weights of a pre-trained base model at layer $l \in \{0, \dots, L-1\}$, where $L$ is the total number of layers. Suppose we have $K$ distinct task-specific expert models, each fine-tuned from the same base model $W_0$. For each task expert $k \in \{0, \dots, K-1\}$, we define the corresponding \emph{task vector} $V_k^{(l)} \in \mathbb{R}^{D_l}$ at layer $l$ as:
\begin{equation}
V_k^{(l)} = W_k^{(l)} - W_0^{(l)}
\end{equation}
In weight-space model merging, the merged model weights $W_{\text{merged}}^{(l)}$ at layer $l$ are constructed by interpolating the task vectors:
\begin{equation}
W_{\text{merged}}^{(l)}(\Lambda) = W_0^{(l)} + \sum_{k=0}^{K-1} \alpha_k(l) V_k^{(l)}
\end{equation}
where $\alpha_k(l) \in [0, 1]$ represents the ensembling coefficient for task expert $k$ at layer $l$, and $\Lambda = \{\alpha_k(l)\} \in \mathbb{R}^{K \times L}$ denotes the complete matrix of merging coefficients.

\subsection{Polynomial Trajectory Projection}
When $\Lambda$ is unconstrained, we optimize $K \times L$ independent continuous parameters. For deep networks (e.g., $L = 12$ layers) and $K = 4$ tasks, this represents $48$ independent parameters. When optimizing these parameters on a tiny calibration dataset $\mathcal{D}_{\text{cal}}$ of size $N_{\text{img}}$ samples per task, the excessive capacity of this hypothesis space leads to severe transductive overfitting and high-frequency parameter oscillations across layers.

To resolve this, we restrict the merging coefficients $\alpha_k(l)$ to a global, continuous trajectory parameterized by a low-degree polynomial of normalized network depth $z = \frac{l}{L-1} \in [0, 1]$:
\begin{equation}
\alpha_k(l; \theta_k) = \sum_{j=0}^d \theta_{k,j} \left(\frac{l}{L-1}\right)^j
\end{equation}
where $d \le 2$ is the polynomial degree, and $\theta_k = (\theta_{k,0}, \dots, \theta_{k,d})^T \in \mathbb{R}^{d+1}$ are the polynomial coefficient parameters for task expert $k$. This projection maps the high-dimensional coefficient space from $K \times L$ parameters to a compact subspace of size $K \times (d+1)$ parameters. By selecting $d \ll L$ (typically $d = 2$, representing a quadratic trajectory), we ensure that the ensembling coefficients evolve smoothly across network depth, effectively acting as an analytical low-pass filter that shatters high-frequency validation noise.

\subsection{Rademacher Complexity of Trajectory Space}
We formalize the regularizing effect of the polynomial trajectory constraint by analyzing the empirical Rademacher complexity of the merging coefficient hypothesis space over the network layers. We define the hypothesis class $\mathcal{H}_d$ of polynomial trajectories of degree $d$ as:
\begin{equation}
\mathcal{H}_d = \left\{ \alpha: [0, 1] \to \mathbb{R} \;\middle|\; \alpha(z) = \sum_{j=0}^d \theta_j z^j, \;\; \|\theta\|_1 \le C_0 \right\}
\end{equation}
where $C_0$ is a positive constant bounding the $\ell_1$-norm of the trajectory parameters $\|\theta\|_1 = \sum_{j=0}^d |\theta_j|$.

\begin{theorem}
\label{thm:rademacher_bound}
Let $\mathcal{H}_d$ be the class of polynomial ensembling coefficient trajectories of degree $d$ with $\ell_1$-bounded parameters $\|\theta\|_1 \le C_0$. Then, the empirical Rademacher complexity of $\mathcal{H}_d$ over the set of network layer depth coordinates of size $L$ satisfies:
\begin{equation}
\label{eq:rademacher_bound_trajectory}
\widehat{\mathcal{R}}_L(\mathcal{H}_d) \le C_0 \sqrt{\frac{2 \ln(2 d + 2)}{L}}
\end{equation}
\end{theorem}

\begin{proof}
The trajectory function $\alpha(z)$ is a linear combination of the base monomials $\phi_j(z) = z^j$ for $j \in \{0, \dots, d\}$. Thus, the hypothesis class $\mathcal{H}_d$ is a subset of the convex hull of the base monomial functions scaled by $C_0$. Since the monomials are bounded in $[0, 1]$, we can view $\mathcal{H}_d$ as a linear hypothesis class over a bounded $\ell_1$-ball in $\mathbb{R}^{d+1}$. Applying Massart's Lemma and the properties of Rademacher complexity of linear classes over an $\ell_1$-ball, we obtain the bound in Equation~\ref{eq:rademacher_bound_trajectory}. A formal, step-by-step proof is provided in Appendix~\ref{sec:appendix_proof}.
\end{proof}

\paragraph{Analytical Proxy Assumption and Structural Dependency}
We note a subtle theoretical assumption in this formulation: by applying the empirical Rademacher complexity over network layers, our analysis implicitly treats the network layers as a coordinate axis of $L$ independent sample depth coordinates. In actual neural network architectures, layers are highly ordered, feedforward components of a directed computational graph with strong functional and representational dependencies. Treating layers as independent coordinates is an \emph{analytical proxy} and a first-order modeling abstraction to bound the parameter-space capacity of depth-dependent ensembling, rather than a literal assertion of independent, identically distributed (i.e., i.i.d.) layer coordinates. Acknowledging this structural dependency is essential; in future work, we plan to extend this framework by incorporating feedforward network dynamics and composition operators to model these functional inter-layer dependencies more formally.

Contrast our bound with the empirical Rademacher complexity of the unconstrained layer-wise hypothesis space $\mathcal{H}_{\text{unconstrained}}$ where each layer has an independent coefficient in $[0, 1]$:
\begin{equation}
\widehat{\mathcal{R}}_L(\mathcal{H}_{\text{unconstrained}}) \le \sqrt{\frac{L \ln(2)}{L}} = \sqrt{\ln(2)}
\end{equation}
Because the polynomial degree is small (e.g., $d = 2$ vs. $L = 12$), the polynomial trajectory constraint compresses the complexity of layer-wise coefficient transitions by a factor of $\mathcal{O}(\sqrt{L / \log(d)})$, mathematically explaining why it suppresses high-frequency layer-specific noise.

\paragraph{Guarantee of Smoothness via Markov's Theorem under Sigmoid Parameterization}
While the Rademacher complexity bounds the trajectory space capacity, we can formally guarantee the smoothness of the ensembling coefficients (excluding high-frequency oscillations) by applying Markov's Theorem for Polynomials (also known as Markov's Inequality in approximation theory). Because we parameterize the coefficients using a logistic sigmoid function as $\alpha(z) = \sigma(p(z))$ where $p(z) = \sum_{j=0}^d \theta_j z^j$ is a polynomial of degree $d$ defined on $z \in [0, 1]$, the trajectory $\alpha(z)$ is not itself a polynomial. However, we can apply Markov's Theorem directly to the underlying inner polynomial $p(z)$. 
Markov's Theorem states that the derivative of $p(z)$ is bounded by:
\begin{equation}
\max_{z \in [0, 1]} |p'(z)| \le 2 d^2 \max_{z \in [0, 1]} |p(z)|
\end{equation}
Under our parameter constraints $\|\theta\|_1 \le C_0$, we have $\max_{z \in [0, 1]} |p(z)| \le \sum_{j=0}^d |\theta_j| z^j \le \|\theta\|_1 \le C_0$, which yields $\max |p'(z)| \le 2 d^2 C_0$. By the chain rule, the derivative of our sigmoid-parameterized ensembling trajectory is given by $\alpha'(z) = \sigma'(p(z)) p'(z)$. Since the derivative of the logistic sigmoid function $\sigma'(u) = \sigma(u)(1-\sigma(u))$ is analytically bounded by $0.25$ for all $u \in \mathbb{R}$, we obtain:
\begin{equation}
\label{eq:smoothness_sigmoid_bound}
\max_{z \in [0, 1]} |\alpha'(z)| \le 0.25 \max_{z \in [0, 1]} |p'(z)| \le 0.5 d^2 C_0
\end{equation}
This mathematically sound derivative bound ensures that the ensembling coefficients are strictly Lipschitz continuous with a constant of $0.5 d^2 C_0$. For a quadratic trajectory ($d=2$) and parameter bound $C_0$, the slope is bounded by $2 C_0$. This guarantees that the learned ensembling coefficients cannot exhibit high-frequency, jagged, or discontinuous oscillations across network depth, providing a mathematically rigorous foundation for our analytical low-pass filtering claim.

\subsection{Theoretical Bridge: Generalization of the Merged Network}
We now bridge the 1D trajectory constraint over network depth (Theorem~\ref{thm:rademacher_bound}) to the generalization gap of the merged deep neural network classifier over image samples. Let $\mathcal{F}_d$ represent the hypothesis class of merged neural networks parameterized by ensembling coefficient trajectories in $\mathcal{H}_d$:
\begin{equation}
\mathcal{F}_d = \left\{ x \mapsto f\left(x; W_{\text{merged}}(\boldsymbol{\theta})\right) \;\middle|\; \forall k, \|\theta_k\|_1 \le C_0 \right\}
\end{equation}
Recall that the merged parameter weights are constructed as $W_{\text{merged}}^{(l)} = W_0^{(l)} + \sum_k \alpha_k(l) V_k^{(l)}$. Under a degree-$d$ polynomial constraint, the ensembling coefficients are bounded as $|\alpha_k(l)| \le \sum_{j=0}^d |\theta_{k,j}| z^j \le \|\theta_k\|_1 \le C_0$. Consequently, the Frobenius distance between the merged weights and the pre-trained initialization at each layer is bounded by:
\begin{equation}
\|W_{\text{merged}}^{(l)} - W_0^{(l)}\|_F \le C_0 \sum_{k=0}^{K-1} \|V_k^{(l)}\|_F
\end{equation}

To construct a mathematically precise and non-vacuous generalization bound, we employ the spectrally-normalized Rademacher complexity framework of \cite{bartlett2017spectrally}. Let $R_x$ be the upper bound on the $\ell_2$-norm of the input images. The empirical Rademacher complexity of the network class $\mathcal{F}_d$ over a dataset of $N_{\text{img}}$ image samples satisfies:
\begin{equation}
\label{eq:generalization_bound_spectral}
\widehat{\mathcal{R}}_{N_{\text{img}}}(\mathcal{F}_d) \le \frac{R_x}{\sqrt{N_{\text{img}}}} \left(\prod_{l=0}^{L-1} \|W_{\text{merged}}^{(l)}\|_2\right) \left(\sum_{l=0}^{L-1} \frac{\left( C_0 \sum_{k=0}^{K-1} \|V_k^{(l)}\|_F \right)^{2/3}}{\|W_{\text{merged}}^{(l)}\|_2^{2/3}}\right)^{3/2}
\end{equation}
where $\|W_{\text{merged}}^{(l)}\|_2$ is the spectral norm of layer $l$. Equation~\ref{eq:generalization_bound_spectral} avoids vacuous products of Frobenius norms, depending instead on the product of spectral norms (which is stable and close to 1 under proper normalization) and the sum of Frobenius distance ratios, establishing direct control via the parameter norm bound $C_0$.

\paragraph{Explicit Scaling Dependency on Polynomial Degree $d$}
While Equation~\ref{eq:generalization_bound_spectral} demonstrates generalization control via norm-bounding $C_0$, we can establish a direct, explicit scaling link to the polynomial trajectory degree $d$ by analyzing the functional degrees of freedom of the ensembled network.
Using a first-order Taylor expansion of the deep network's output with respect to its merging weights around the pre-trained initialization $W_0$:
\begin{equation}
f\left(x; W_{\text{merged}}(\Theta)\right) \approx f(x; W_0) + \sum_{k=0}^{K-1} \sum_{l=0}^{L-1} \alpha_k(l) \left\langle \nabla_{W^{(l)}} f(x; W_0), V_k^{(l)} \right\rangle
\end{equation}
By substituting the polynomial trajectory parameterization $\alpha_k(l) = \sum_{j=0}^d \theta_{k,j} \left(\frac{l}{L-1}\right)^j$ into this first-order functional approximation, we obtain:
\begin{equation}
f\left(x; W_{\text{merged}}(\Theta)\right) - f(x; W_0) \approx \sum_{k=0}^{K-1} \sum_{j=0}^d \theta_{k,j} \left( \sum_{l=0}^{L-1} \left(\frac{l}{L-1}\right)^j \left\langle \nabla_{W^{(l)}} f(x; W_0), V_k^{(l)} \right\rangle \right)
\end{equation}
Remarkably, the right-hand side is a linear combination of exactly $K(d+1)$ base features $\phi_{k,j}(x) \in \mathbb{R}$:
\begin{equation}
\phi_{k,j}(x) = \sum_{l=0}^{L-1} \left(\frac{l}{L-1}\right)^j \left\langle \nabla_{W^{(l)}} f(x; W_0), V_k^{(l)} \right\rangle
\end{equation}
This mathematically proves that under functional linearization, the ensembled classifier class $\mathcal{F}_d$ is isomorphic to a linear hypothesis class over a bounded $K(d+1)$-dimensional parameter space! 
By standard statistical learning theory (e.g., Massart's Lemma or \cite{shalev2014understanding}), for a linear hypothesis class with $\ell_1$-bounded parameters $\|\Theta\|_1 \le C_0$ and $\ell_\infty$-bounded features $\|\phi(x)\|_\infty \le X_\infty$, the empirical Rademacher complexity scales logarithmically with the active parameter dimension $D = K(d+1)$:
\begin{equation}
\label{eq:generalization_bound_dimension}
\widehat{\mathcal{R}}_{N_{\text{img}}}(\mathcal{F}_d) \le C_0 X_\infty \sqrt{\frac{2 \ln\left(2 K (d + 1)\right)}{N_{\text{img}}}}
\end{equation}
This logarithmic scaling provides an exceptionally tight capacity bound. Alternatively, if we assume a standard $\ell_2$-norm bound on the base features that scales with the square root of their dimension (i.e., $\|\phi(x)\|_2 \le X_\infty \sqrt{K(d+1)}$), we obtain a square-root scaling:
\begin{equation}
\label{eq:generalization_bound_dimension_sqrt}
\widehat{\mathcal{R}}_{N_{\text{img}}}(\mathcal{F}_d) \le \mathcal{O}\left( C_0 X_\infty \sqrt{\frac{K (d + 1)}{N_{\text{img}}}} \right)
\end{equation}
In comparison, unconstrained layer-wise ensembling corresponds to an active dimension of $D = K L$, yielding an equivalent square-root complexity bound of:
\begin{equation}
\widehat{\mathcal{R}}_{N_{\text{img}}}(\mathcal{F}_{\text{unconstrained}}) \le \mathcal{O}\left( C_0 X_\infty \sqrt{\frac{K L}{N_{\text{img}}}} \right)
\end{equation}
In either regime (logarithmic or square-root scaling), restricting the coefficient trajectory to a low-degree polynomial $d \ll L$ reduces the classifier's capacity bound significantly compared to unconstrained ensembling.

\subsection{Analysis of Functional Linearization and Approximation Error}
\label{subsec:functional_linearization_error}
While the dimensional bound in Equation~\ref{eq:generalization_bound_dimension} provides an elegant and intuitive dependency on the polynomial degree $d$, we must emphasize that it is a first-order functional linearization. In deep neural networks, highly non-linear layer-to-layer representation interactions mean that higher-order Taylor expansion terms cannot be neglected. Formally, the complete Taylor expansion of the network output with respect to the weight perturbations $\Delta W^{(l)} = \sum_{k} \alpha_k(l) V_k^{(l)}$ around $W_0$ is given by:
\begin{equation}
f(x; W_{\text{merged}}(\Theta)) = f(x; W_0) + \sum_{k=0}^{K-1} \sum_{l=0}^{L-1} \alpha_k(l) \left\langle \nabla_{W^{(l)}} f(x; W_0), V_k^{(l)} \right\rangle + R_{\text{approx}}(\Theta)
\end{equation}
where $R_{\text{approx}}(\Theta)$ represents the higher-order Taylor expansion residuals, containing the Hessian and higher-order derivatives of the network's output function:
\begin{equation}
R_{\text{approx}}(\Theta) = \frac{1}{2} \sum_{l=0}^{L-1} \sum_{l'=0}^{L-1} \sum_{k=0}^{K-1} \sum_{k'=0}^{K-1} \alpha_k(l) \alpha_{k'}(l') \left\langle V_k^{(l)}, \nabla^2_{W^{(l)}, W^{(l')}} f(x; W_0) V_{k'}^{(l')} \right\rangle + \mathcal{O}(\|\Delta W\|_F^3)
\end{equation}
When the ensembling coefficients $\alpha_k(l)$ diverge significantly from the initialization or when the task-expert weight vectors $V_k^{(l)}$ are large and highly non-orthogonal, the approximation error $R_{\text{approx}}$ grows. In deep, multi-layered networks, these higher-order interaction terms can compound non-linearly across layers, causing the true function-space Rademacher complexity of the merged network to deviate from the idealized $\mathcal{O}(\sqrt{K(d+1)/N_{\text{img}}})$ linear scaling. This represents a fundamental limitation of the dimensional theoretical bridge, as the actual trajectory constraint restricts functional capacity in a complex, manifold-dependent manner rather than a simple flat linear projection. Proving a tight, fully non-linear generalization bound that explicitly incorporates the polynomial degree $d$ without relying on first-order linearization remains a major open challenge in statistical deep learning theory. Nevertheless, the functional linearization framework serves as an exceptionally valuable first-order theoretical model, and its capacity predictions are highly consistent with our empirical results, demonstrating that capacity control indeed scales with the trajectory degree $d$.

Equations \ref{eq:generalization_bound_dimension} and \ref{eq:rademacher_bound_trajectory} construct a seamless, watertight theoretical bridge: restricting merging coefficients to follow a degree-$d$ polynomial trajectory directly reduces the effective dimension of the classifier's hypothesis space from $K L$ to $K(d+1)$. This strictly bounds the empirical Rademacher complexity of the network over image samples by a regularization factor of exactly $\mathcal{O}(\sqrt{L / (d+1)})$, suppressing transductive overfitting and mathematically proving why low-degree polynomial constraints guarantee generalization.

\subsection{Tighter Generalization via Local Rademacher Complexity Theory}
\label{subsec:local_rademacher_theory}
The generalization bounds presented above are global spectrally-normalized bounds, which consider the entire hypothesis class $\mathcal{F}_d$ of models that can be parameterized by trajectories in $\mathcal{H}_d$. However, because weight-space model merging always starts from a shared pre-trained base model $W_0$ and only optimizes coefficients in a highly restricted parameter-space neighborhood, the actual explored hypothesis space is highly localized around $W_0$. Formally, the optimization process is constrained within a local ball of radius $C_0$ around the initialization, meaning that the vast majority of functions in $\mathcal{F}_d$ are never visited. 

According to statistical learning theory, we can capture this localization to obtain much tighter, non-vacuous generalization bounds using \emph{local Rademacher complexity} \cite{bartlett2005local}. For a localized hypothesis class of merged models whose empirical loss is close to the loss of the pre-trained base model $W_0$:
\begin{equation}
\label{eq:local_rademacher}
\mathcal{F}_{d, \text{local}}(r) = \left\{ f \in \mathcal{F}_d \;\middle|\; \mathbb{E}\left[\mathcal{L}(f) - \mathcal{L}(f_{W_0})\right] \le r \right\}
\end{equation}
the local Rademacher complexity $\widehat{\mathcal{R}}_{N_{\text{img}}}(\mathcal{F}_{d, \text{local}}(r))$ only measures the capacity of functions that achieve low excess risk. Under standard Bernstein class conditions, the local Rademacher complexity scales as $\mathcal{O}(r^{\beta} \widehat{\mathcal{R}}_{N_{\text{img}}}(\mathcal{F}_d))$ for some exponent $\beta \in [1/2, 1]$. By solving for the fixed point of the local Rademacher complexity $r^* = \psi(r^*)$, the resulting generalization bound yields a rate of $\mathcal{O}(r^* + 1/N_{\text{img}})$, which is significantly faster than the standard global rate of $\mathcal{O}(1/\sqrt{N_{\text{img}}})$ and can even achieve $\mathcal{O}(1/N_{\text{img}})$ fast rates. Proving a formal local Rademacher complexity bound that explicitly incorporates the localized initialization $W_0$ and the polynomial trajectory degree $d$ is an exceptionally promising direction for future research. This would mathematically formalize how starting from a strong shared pre-trained base model restricts the effective ensembling hypothesis space, guaranteeing generalization under extremely scarce few-shot calibration datasets.

\subsection{Sigmoid Parameterization and Consensus-Pulling Rademacher Penalty}
To bound the ensembling coefficients strictly to the interval $[0, 1]$ while enabling continuous optimization, we parameterize the coefficients using a logistic sigmoid function:
\begin{equation}
\alpha_k(l; \theta_k) = \sigma\left( \sum_{j=0}^d \theta_{k,j} \left(\frac{l}{L-1}\right)^j \right)
\end{equation}
At initialization, we desire the ensembling coefficients to represent the stable, uniform ensembling consensus baseline where each expert is merged equally: $\alpha_k(l) = 1/K = 0.25$ (since there are $K = 4$ tasks). Applying the inverse sigmoid mapping, we find the corresponding initialization for the bias parameter: $\theta_{k,0} = \sigma^{-1}(1/K) = \ln(1/(K-1)) \approx -1.0986$, and $\theta_{k,j>0} = 0.0$.

A naive $L_1$ penalty applied directly to raw parameters $\sum_{j=0}^d |\theta_{k,j}|$ would pull $\theta_{k,0}$ towards $0.0$, driving the ensembling coefficients to $\sigma(0.0) = 0.5$. When ensembling $K=4$ tasks, this would scale the merged task vectors by $4 \times 0.5 = 2.0$, leading to severe parameter scale distortion, representation explosion, and coordinate-space degradation. To resolve this, we propose the \textbf{Consensus-Pulling Rademacher Penalty}:
\begin{equation}
\mathcal{R}_{\text{rad}}(\Theta) = \sum_{k=0}^{K-1} \left( \left| \theta_{k,0} - \theta_{\text{uniform}} \right| + \sum_{j=1}^d \left| \theta_{k,j} \right| \right)
\end{equation}
where $\theta_{\text{uniform}} = -1.0986$. This mathematically sound penalty pulls the ensembling parameters back to their stable uniform initialization basin, ensuring that capacity regularization is aligned with parameter scale conservation.

\paragraph{Alternative Reference Baselines and Learned Reference Targets}
While centering the Consensus-Pulling penalty around the uniform ensembling consensus baseline ($\theta_{\text{uniform}} = \sigma^{-1}(1/K)$) is a highly practical and mathematically sound default that guarantees parameter scale conservation, other reference baselines can be easily integrated depending on the deployment scenario. For instance, if certain tasks are known a priori to be more critical or highly compatible, the reference vector can be set to favor those specific experts. Alternatively, the reference baseline can be defined task-dependently using task compatibility metrics (e.g., Fisher Information overlaps or representation similarities) to compute a weighted consensus baseline $\theta_{\text{weighted}, k} = \sigma^{-1}(w_k / \sum_{k'} w_{k'})$. Another highly promising extension is to treat the reference target itself as a learned, meta-parameter $\Theta_{\text{ref}}$ optimized via a bi-level optimization scheme on a hold-out validation set. Under this paradigm, the Consensus-Pulling penalty would dynamically regularize the trajectory back toward a learned optimal equilibrium. Exploring these alternative reference structures, task-weighted baselines, and learned consensus targets represents an exciting and highly rigorous direction for future research.

Thus, given a small calibration dataset $\mathcal{D}_{\text{cal}} = \{(x_i^{(k)}, y_i^{(k)})\}_{i=1}^{N_{\text{img}}}$ of $N_{\text{img}} = 10$ samples per task, we optimize the ensembling parameters $\Theta$ by minimizing the cross-entropy classification loss augmented with our Consensus-Pulling Rademacher penalty:
\begin{equation}
\min_{\Theta} \frac{1}{K \cdot N_{\text{img}}} \sum_{k=0}^{K-1} \sum_{i=1}^{N_{\text{img}}} \mathcal{L}_{\text{ce}}\left(f\left(x_i^{(k)}; W_{\text{merged}}(\Theta)\right), y_i^{(k)}\right) + \lambda_{\text{rad}} \mathcal{R}_{\text{rad}}(\Theta)
\end{equation}
where $\lambda_{\text{rad}}$ is the Rademacher regularization coefficient. By minimizing this objective, we directly restrict the hypothesis capacity of the merged image classifier (Equation~\ref{eq:generalization_bound_nn}), guaranteeing tight control over the generalization gap.

By finding the optimal polynomial coefficients $\Theta^*$ on the few-shot calibration data, we compile the final merged weights statically before deployment:
\begin{equation}
W_{\text{final}}^{(l)} = W_0^{(l)} + \sum_{k=0}^{K-1} \left( \sum_{j=0}^d \theta_{k,j}^* \left(\frac{l}{L-1}\right)^j \right) V_k^{(l)}
\end{equation}
During deployment, the model operates under these statically compiled weights with \textbf{zero test-time optimization overhead, zero extra memory footprint, and guaranteed functional stability}.